\begin{enunciado}{\ejercicio}
Hallar, si existe, una matriz $\boldsymbol{A}$ con coeficientes reales y del tamaño adecuado para que los valores singulares no nulos de $\boldsymbol{A}$ sean $\left\{\frac{3}{2}, 3\right\}$,
$$
\boldsymbol{A} \begin{pmatrix} 1 \\ -2 \\ 2 \end{pmatrix} = \begin{pmatrix} 0 \\ 0 \\ 0 \end{pmatrix} \quad \text{y} \quad \begin{pmatrix} 2 & 2 & 1 \end{pmatrix} \boldsymbol{A} = \begin{pmatrix} 0 & 0 & 0 \end{pmatrix}.
$$
\end{enunciado}

Para construir la matriz $A$ la idea es aprovechar propiedades que nos da la descomposición SDV de $A = U \Sigma V^t$, primero notemos qué información nos da el enunciado:

Sabemos que $\{\frac{3}{2}, 3\}$ son valores singulares de A, por lo tanto, estos van a ser algunos de los $\sigma_i$ de $\Sigma$.

El vector $\begin{pmatrix} 1 \\ -2 \\ 2 \end{pmatrix}$ pertenece al núcleo de $A$, lo que quiere decir que este vector es alguno de los $\{v_{r+1}, \dots, v_m\} \in Nu(A)$.

El vector $\begin{pmatrix} 2 \\ 2 \\ 1 \end{pmatrix}$ pertenece al núcleo de $A^t$, esto se debe a que si hacemos algunas cuentas nos queda que:

$$
\begin{pmatrix} (2 , 2, 1 \end{pmatrix} A)^t = \begin{pmatrix} 0,0,0 \end{pmatrix} ^t
$$
$$
A^t \begin{pmatrix} 2 , 2, 1 \end{pmatrix}^t = \begin{pmatrix} 0 \\ 0 \\ 0 \end{pmatrix} \\
$$
$$
A^t \begin{pmatrix} 2 \\ 2 \\ 1 \end{pmatrix} = \begin{pmatrix} 0 \\ 0 \\ 0 \end{pmatrix} \\
$$


si dicho vector pertenece al núcleo de $A^t$ entonces es alguno del conjunto: $\{u_{r+1}, \dots, u_n\} \in Nu(A^t)$

Ahora, usando esta información veamos que dimensión puede tener $A$:
$$
    A = \ua{U}{m\times m} \ \oa{\Sigma}{m \times n} \ \ua{V^t}{n \times n}
$$
por lo dicho antes, sabemos que $\begin{pmatrix} 2 \\ 2 \\ 1 \end{pmatrix} \in U$ y que $\begin{pmatrix} 1 , -2, 2 \end{pmatrix} \in V$, y como las dimensiones de $U$ y $V$ deben de ser cuadradas, seguro que $m = 3$ y $n =3$ funcionaría para encontrar una posible A, entonces la SVD nos quedaría más o menos así: (observemos que normalizamos los vectores)

$$
A = U \Sigma V^*
= \begin{pmatrix}
    \mid & \mid & \frac{2}{3} \\
    u_1 & u_2 & \frac{2}{3} \\
    \mid & \mid & \frac{1}{3}
    \end{pmatrix}
    \begin{pmatrix}
    3 & 0 & 0 \\
    0 & \frac{3}{2} & 0 \\
    0 & 0 & 0
    \end{pmatrix}
    \begin{pmatrix}
    \text{---} & v_1^* & \text{---} \\
    \text{---} & v_2^* & \text{---} \\
    \frac{1}{3} & -\frac{2}{3} & \frac{2}{3}
    \end{pmatrix}
$$

Notemos además que $\sigma_3$ debe de ser 0, si no, estarían en la imagen los vectores $u_3$ y $v_3*$ que dijimos que estaban en el núcleo.

Ahora debemos buscar los vectores que nos faltan, arranquemos por completar $U$:

Sabemos que $u_1, u_2, u_3$ (con $u_3 = (\frac{2}{3}, \frac{2}{3}, \frac{1}{3})$) debe de formar una base ortonormal para que $U$ sea unitaria.

construyámosla: se puede hacer a ojo o usando Gram-Schmidt. A ojo saco a $u_1$:
$$
u_1 = \begin{pmatrix} \frac{1}{\sqrt{2}} \\ -\frac{1}{\sqrt{2}} \\ 0\end{pmatrix}
$$
El vector $u_2$ no lo saco a ojo, usemos GS para hallarlo aplicándolo a un canónico cualquiera:
$$
\begin{align*}
u_3 &= (0,0,1) - \text{proy}_{(0,0,1)}(u_1) - \text{proy}_{(0,0,1)}(u_3) \\
    &= (0,0,1) - 0 \cdot u_1 - \frac{1}{3} (\frac{2}{3}, \frac{2}{3}, \frac{1}{3}) \\
    &= (0,0,1) - (\frac{2}{9}, \frac{2}{9}, \frac{1}{9}) \\
    &= (-\frac{2}{9}, -\frac{2}{9}, \frac{8}{9})
\end{align*}
$$
Normalizándolo, nos queda:
$$
u_3 = \frac{u_3}{||u_3||} = \frac{(-\frac{2}{9}, -\frac{2}{9}, \frac{8}{9})}{ \frac{2 \sqrt{2}}{3}} = \frac{3}{2 \sqrt{2}} (-\frac{2}{9}, -\frac{2}{9}, \frac{8}{9}) = (-\frac{1}{3\sqrt{2}},-\frac{1}{3\sqrt{2}},\frac{4}{3\sqrt{2}})
$$
Entonces, por ahora nuestra matriz A va quedando de la siguiente manera:
$$
A = U \Sigma V^*
= \begin{pmatrix}
    \frac{1}{\sqrt{2}} & -\frac{1}{3 \sqrt{2}} & \frac{2}{3} \\
    -\frac{1}{\sqrt{2}} & -\frac{1}{3 \sqrt{2}} & \frac{2}{3} \\
    0 & \frac{4}{3 \sqrt{2}} & \frac{1}{3}
    \end{pmatrix}
    \begin{pmatrix}
    3 & 0 & 0 \\
    0 & \frac{3}{2} & 0  \\
    0 & 0 & 0
    \end{pmatrix}
    \begin{pmatrix}
    \text{---} & v_1^* & \text{---} \\
    \text{---} & v_2^* & \text{---} \\
    \frac{1}{3} & -\frac{2}{3} & \frac{2}{3}
    \end{pmatrix}
$$

Debemos buscar los $v_i$ expandiendo $\{v_1*, v_2*, v_3*\}$ (con $v_3* = (\frac{1}{3} , -\frac{2}{3} , \frac{2}{3} )$) a una base ortonormal (como hicimos con $U$). Por suerte, los vectores $u_3$ y $v_3$ són bastante parecidos, así que elegimos los mismos vectores que antes, pero cambiando los signos necesarios asi $V$ queda unitaria ;), nos queda $V$ como:

$$
V = \begin{pmatrix}
\frac{4}{3 \sqrt{2}} & \frac{1}{3 \sqrt{2}} & -\frac{1}{3 \sqrt{2}} \\
0 & -\frac{1}{\sqrt{2}} & -\frac{1}{\sqrt{2}} \\
\frac{1}{3} & - \frac{2}{3} & \frac{2}{3}
\end{pmatrix}
$$

Ahora, nuestra A quedaría de la siguiente manera:

$$
A = U \Sigma V^*
    = \begin{pmatrix}
        \frac{1}{\sqrt{2}} & -\frac{1}{3 \sqrt{2}} & \frac{2}{3} \\
        \frac{-1}{\sqrt{2}} & -\frac{1}{3 \sqrt{2}} & \frac{2}{3} \\
        0 & \frac{4}{3 \sqrt{2}} & \frac{1}{3}
    \end{pmatrix}
    \begin{pmatrix}
        3 & 0 & 0 \\
        0 & \frac{3}{2} & 0  \\
        0 & 0 & 0
    \end{pmatrix}
    \begin{pmatrix}
        \frac{4}{3 \sqrt{2}} & \frac{1}{3 \sqrt{2}} & -\frac{1}{3 \sqrt{2}} \\
        0 & -\frac{1}{\sqrt{2}} & -\frac{1}{\sqrt{2}} \\
        \frac{1}{3} & - \frac{2}{3} & \frac{2}{3}
    \end{pmatrix}
$$
$$
    = \begin{pmatrix}
            \frac{1}{\sqrt{2}} & -\frac{1}{3 \sqrt{2}} & \frac{2}{3} \\
            \frac{-1}{\sqrt{2}} & -\frac{1}{3 \sqrt{2}} & \frac{2}{3} \\
            0 & -\frac{4}{3 \sqrt{2}} & \frac{1}{3}

            \end{pmatrix}
            \begin{pmatrix}
            2 \sqrt{2} & \frac{\sqrt{2}}{2} & -\frac{\sqrt{2}}{2} \\
            0 & -\frac{3 \sqrt{2}}{4} & - \frac{3 \sqrt{2}}{4}\\
            0 & 0 & 0
            \end{pmatrix}
    = \begin{pmatrix}
                2 & \frac{3}{4} & -\frac{1}{4} \\
                -2 & - \frac{1}{4} & \frac{3}{4}  \\
                0 & -1 & -1
        \end{pmatrix}
$$

Podemos concluir que, una posible A es:
$$
A = \begin{pmatrix}
                2 & \frac{3}{4} & -\frac{1}{4} \\
                -2 & - \frac{1}{4} & \frac{3}{4}  \\
                0 & -1 & -1
        \end{pmatrix}
$$

\begin{aportes}
  \item \aporte{\dirRepo}{Luloide \github}
\end{aportes}
